\clearpage
\section*{ÖZET}
\label{ozet}

%\par Projeyi kısaca anlatıp, çok fazla ayrıntıya girmeden okuyacak kişiye konu ve çalışmanın içeriği ile ilgili bilgi veren, 100 kelimeyi aşmayacak şekilde çalışmanızı özetleyiniz.
%\par Problemden, ugulamış olduğunuz çözümden, elde ettiğiniz sonuçalrdan kısaca bahsediniz.

%Buraya giriş için bir cümle eklenebilir.

\par Bu çalışmada, model roket sistemlerinde deterministik, modüler ve genişletilebilir bir uçuş kontrol yazılımının oluşturulması hedeflenmiştir. Bu doğrultuda, STM32F446 mikrodenetleyicisi üzerinde FreeRTOS tabanlı bir uçuş kontrol bilgisayarı mimarisi tasarlanmıştır. Çalışma sürecinde, roketlerde kullanılan gerçek zamanlı işletim sistemlerine yönelik kapsamlı bir literatür taraması gerçekleştirilmiştir. Ayrıca, kullanılan sensörlerin ve modüllerin teknik dokümanları incelenmiştir.

\par Bitirme Projesi II kapsamında, birinci aşamada oluşturulan temel yazılım altyapısının üzerine doğrudan uçuşa yönelik kritik kontrol algoritmaları ve gelişmiş sensör füzyon yöntemleri entegre edilmiştir. Bu doğrultuda, BMI088 ve BME280 sensörlerinden kesme ve DMA tabanlı veri okuma mekanizmaları hayata geçirilerek CPU tıkanıklığı önlenmiştir. Durum kestirimi için Mahony AHRS ve Yükseklik Kalman filtresi entegre edilerek, uçuş fazlarını milisaniyeler seviyesinde yöneten 7-fazlı bir Uçuş Durum Makinesi (flight\_sm) oluşturulmuştur. Geliştirilen algoritmalar ve donanımsal arıza önleme mekanizmaları (IWDG), Processor-in-the-Loop (PIL) test metodolojisi ve RocketPy simülasyonları ile doğrulanmış, gürültülü sensör verilerinde dahi hatasız ve kararlı çalışan tam işlevsel bir aviyonik uçuş kontrol yazılımı elde edilmiştir.


\vspace{3mm}

 \begin{tabular}{ll}
 	{\textbf{Anahtar Kelimeler}:} & {Model Roket, Uçuş Bilgisayarı, FreeRTOS, Kalman Filtresi, Sensör Füzyonu}  \\   	
 \end{tabular} 




 


